\chapter{Materials and Methods}
\section{Study Design}
This thesis uses a computational retrospective study design based on publicly available scRNA-seq data. The workflow begins with raw count matrices and metadata, continues through preprocessing and feature selection, and ends with supervised classification and biomarker ranking.

\section{Dataset}
The dataset was obtained from GSE183276, the Human Kidney Healthy and Injury Cell Atlas. It includes human kidney cells profiled with 10x Chromium scRNA-seq technology and annotated by disease condition. The primary class label is \texttt{condition.l1}, containing \texttt{Ref}, \texttt{AKI}, and \texttt{CKD}.
\begin{table}[H]
\centering
\caption{Summary of the GSE183276 dataset used in this thesis.}
\label{tab:dataset_summary}
\begin{tabularx}{\textwidth}{lX}
\toprule
Property & Description \\
\midrule
Dataset accession & GSE183276 \\
Dataset title & Human Kidney Healthy and Injury Cell Atlas \\
Organism & \textit{Homo sapiens} \\
Technology & 10x Chromium single-cell RNA sequencing \\
Input files & \texttt{counts.mtx}, \texttt{genes.tsv}, \texttt{barcodes.tsv}, metadata file \\
Processed file & \texttt{GSE183276\_HVG3000\_by\_variance.h5ad} \\
Primary label & \texttt{condition.l1} \\
Classes & Ref, AKI, CKD \\
Feature selection & Top 3000 highly variable genes by variance \\
\bottomrule
\end{tabularx}
\end{table}


\section{Computational Environment}
The analysis was implemented in Python using Scanpy, AnnData, NumPy, pandas, SciPy, scikit-learn, Matplotlib, and Seaborn. LaTeX was used for thesis preparation, with IEEE citation style for references.

\section{Overall Workflow}
\begin{figure}[H]
\centering
\begin{tikzpicture}[node distance=0.95cm, every node/.style={font=\small}, box/.style={rectangle, rounded corners, draw=black, fill=blue!7, minimum width=5.8cm, minimum height=0.8cm}, arr/.style={-{Latex}, thick}]
\node[box] (a) {Raw scRNA-seq counts and metadata};
\node[box, below=of a] (b) {Quality control filtering};
\node[box, below=of b] (c) {Normalization and log transformation};
\node[box, below=of c] (d) {Top 3000 highly variable genes};
\node[box, below=of d] (e) {Random Forest classification};
\node[box, below=of e] (f) {Pairwise biomarker discovery};
\node[box, below=of f] (g) {Biological interpretation};
\draw[arr] (a)--(b); \draw[arr] (b)--(c); \draw[arr] (c)--(d); \draw[arr] (d)--(e); \draw[arr] (e)--(f); \draw[arr] (f)--(g);
\end{tikzpicture}
\caption{Overall computational workflow for CKD and AKI biomarker discovery.}
\label{fig:overall_workflow}
\end{figure}
